\subsection*{The MaxToki-217M case study}
This part reports what the corrected analyses show about MaxToki-217M.
Every result concerns MaxToki-217M given one cell at a time, with the
input encoding and the intervention code checked as described in
Materials and methods. Table~\ref{tab:scope} lists, for each analysis,
the cells it used, the reference it was judged against, and what its
conclusion depends on. Four deployed analyses used incorrectly encoded
inputs and were not re-run, so we do not report them: the topology
screen of gene--gene representations, the ageing probe, the
feature-patching specificity ratio, and the edge census from 1,000
layer-5 features in 20 cells, which also used the intervention code that
deleted a transformer block. The MaxToki-1B attention run and the
Adamson attention run are not reported for the same reason. The
gene-set annotation of the sparse-autoencoder features was redone only
to choose the source features for circuit tracing, and its rates are not
reported. S3~Appendix lists smaller follow-up analyses that were not
re-run either, such as value-weighted attention, head ablation and the
later uses of the circuit edge list.

% TABLE: scope (tab:scope) is inserted here.

\subsubsection*{Attention scores do not beat gene-level baselines}
MaxToki is a causal decoder: genes enter in order of falling normalised
expression, and a gene can attend only to genes placed before it. The
attention from a regulator $P$ to a target $T$ is therefore zero in every
cell in which $T$ comes after $P$, and the average over cells partly
measures how often $T$ is ranked above $P$ (the order share), which
needs no model. We read five attention scores at layer~8: regulator to target, target to
regulator, the mean and the larger of the two directions, and regulator
to target averaged only over the cells in which the target comes before
the regulator. A gene counted as responding to a knockdown when its mean
$\log(1+x)$ value in the knocked-down cells differed from that in 5,000
control cells by at least 0.5 (natural log), with BH $q < 0.05$ over the
1,500 genes tested (Welch $t$-test). Here $x$ is the count scaled so that
each cell sums to 10,000 over those 1,500 genes. A knockdown was kept
when it had at least 3 responding and 3 non-responding genes. The
deployed RPE1 pipeline used this rule. We kept it for K562 and chose it
before seeing the K562 results. On the RPE1 screen, which was encoded
correctly from raw counts, the order share alone predicted knockdown responses better than
the regulator-to-target score (mean AUROC over 325 knockdowns 0.648
against 0.603), and the variance of the target gene predicted them far
better (0.766, 95\% bootstrap interval over knockdowns 0.757--0.776)
(Fig~\ref{fig:caseReg}A). AUROC is the probability that a randomly chosen
responding gene is scored above a randomly chosen non-responding one;
0.5 is chance.

\begin{figure}[H]
\caption{\textbf{MaxToki-217M: regulatory probes after correction.}
(A)~Prediction of which genes respond to a CRISPR-interference knockdown,
in RPE1 cells (325 knockdowns) and K562 cells (194 knockdowns): mean
AUROC over knockdowns (line, 95\% bootstrap interval over knockdowns) for
layer-8 attention read from regulator to target and in both directions,
and for three scores that use no model: the share of cells in which the
target is ranked above the regulator, co-expression, and the variance of
the target gene. Dashed line, chance. (B)~Recovery of TRRUST
regulator--target pairs (RPE1: 16 factors, 138 pairs; K562: 16 factors,
107 pairs): pooled AUROC with 95\% bootstrap interval over factors; the
target-gene model uses each target's mean, variance and dropout rate.
(C)~Power of the count-matched test for GATA1 in K562 cells, counting
every overlapping target (the deployed analysis stopped at 5): share of
200 simulations in which GATA1 ChIP-seq targets planted among the top 20
genes of one responding feature were detected at $p<0.05$ (band, Wilson
95\% interval). (D)~Prediction of the direction of change of 796,304
knockdown--target records in K562 cells (248 silenced genes) by the
traced circuit (line, 95\% bootstrap interval over silenced genes), by
always predicting ``down'', and by predicting each target's usual
direction from other silenced genes.}
\label{fig:caseReg}
\end{figure}

On TRRUST regulator--target pairs in RPE1 (16 transcription factors, 138
pairs), a symmetric attention score (the mean of both directions) reached
an AUROC of 0.638 and beat a degree-preserving network null ($z = +3.2$;
empirical $p = 0.003$ from 1,000 rewired networks)
(Fig~\ref{fig:caseReg}B). After Benjamini--Hochberg (BH) adjustment the
$q$ value was 0.018. The adjusted family had 12 tests: six scores (the
five attention scores and the order share) in each of the two cell
lines. Co-expression computed from the same cells beat the null about as
strongly (AUROC 0.596, $z = +3.0$). Attention scored 0.042 higher, but
the 95\% interval of this gap ($-0.040$ to $+0.127$) includes zero. A
three-number model of each target gene knows nothing about the
regulator. It uses the gene's mean, its variance and its dropout rate
(the share of control cells in which the gene has zero counts). It scored
0.749 (0.691--0.801). The symmetric attention score was 0.11 below this
gene model (0.00 to 0.20). We also removed the mean, variance and
dropout rate of both genes from the symmetric score with the planned
linear model. What remained had an AUROC of 0.538 (0.456--0.614), which
does not differ from 0.5. It still beat the rewired networks
($z = +2.75$, $p = 0.003$, not adjusted). This is because the rewired
networks were scored after the same removal, and their mean AUROC fell
too, from about 0.58 to about 0.48. Co-expression treated the same way
also beat them ($z = +2.94$), and the symmetric score still beat them
after co-expression was removed
(0.611, $z = +2.53$). So in RPE1, symmetric attention carries a
pair-level signal about as large as that of co-expression, not a larger
one (S3~Appendix).

In K562 cells (2,000 control cells of the same screen, 1,500 genes,
layer~8), attention predicted knockdown responses at chance under this
rule. Over 194 knockdowns the mean AUROC was 0.505 (95\%
bootstrap interval over knockdowns 0.491--0.520) from regulator to target
and 0.502 (0.488--0.517) for the symmetric score; all five attention
scores lay between 0.499 and 0.509 (Fig~\ref{fig:caseReg}A). Target
variance scored 0.690 (0.674--0.708) and co-expression 0.556
(0.540--0.573). A second label rule ran the same test on the stored
$\log(1+\text{CP10k})$ values. These are scaled over all 6,546 genes in
the file, not over the 1,500 genes. Under this rule only 67 knockdowns
had at least 3 responding and 3 non-responding genes. With it
the regulator-to-target score was 0.550 (0.517--0.580). That is above
chance, but below the order share (0.611, 0.568--0.651), which needs no
model, and below target variance (0.695, 0.659--0.730) (S3~Appendix).
On the 107 TRRUST pairs of 16 factors in this gene set, the symmetric
score had a pooled AUROC of 0.571 (0.521--0.617) and $z = +1.76$ against
the rewired networks (empirical $p = 0.041$; BH $q = 0.069$ over the 12
tests), so the RPE1 result did not replicate in K562
(Fig~\ref{fig:caseReg}B). The regulator-to-target score had a pooled
AUROC of 0.548 (0.499--0.594) and $z = +1.72$ ($q = 0.069$), and it
rested on one factor. MYC holds 27 of the 107 pairs, and without MYC
this $z$ fell to $+0.12$ (and the symmetric score's $z$ to $+1.11$).
Co-expression from the same cells did as well as attention (0.588,
$z = +1.55$), and the target-gene model did better (0.723,
0.644--0.814). After the same gene-level features were removed, no K562
score beat the rewired networks (lowest BH $q = 0.18$ over the six
scores with the planned linear model, and 0.17 with a boosted-tree
model). Which genes enter the test matters. Our gene set is the 1,500
genes with the highest variance of log-normalised expression. The
deployed analysis used the same rule on values that had been logged a
second time, and only 627 of its 1,500 genes are in our set. On the
deployed gene set (8 factors, 45 pairs), the symmetric score did beat
the null ($z = +2.38$, $p = 0.010$). This check was not planned and is
not adjusted. Single heads, other layers as a primary endpoint and
references other than TRRUST were not tested (S3~Appendix).

\subsubsection*{No sparse-autoencoder feature is specific to a knocked-down transcription factor}
A sparse autoencoder (SAE) rewrites a hidden state as the sum of a few
learned directions, called features, out of a much larger set. We
trained twelve SAEs on 500 K562 control cells, one for each hidden state
the model exposes. Layer~0 is the token embedding. Layers 1--10 are the
outputs of the first ten of the 11 transformer blocks; this hidden state,
passed from one block to the next, is called the residual stream.
Layer~11 is the output of the last block after the final normalisation.
Each SAE has 4,928 features and keeps only the 32 most active features
for each token (a TopK SAE). The SAEs explained 70\% of the held-out
variance at layer~0 and 84--93\% at layers 1--11. For 87 transcription
factors knocked down in K562 (34 to 100 cells each; 20 with ChIP-seq
target sets in DoRothEA~\cite{garcia-alonso2019dorothea}), we called a
layer-5 feature responsive when its activation differed from control
cells (Mann--Whitney test, BH $q < 0.05$, and a minimum effect size),
and asked whether the 20 genes on which a responsive feature is most
active overlap the factor's targets more than expected. No factor passed
after BH correction across factors, at any of seven effect cut-offs and
under any of four nulls. In the count-matched null, each top gene is
swapped for a random gene seen in a similar number of cells. The second
null also matches gene length. The third uses random groups of control
cells that keep as many responding features as the knockdown. The fourth
uses the responding features of other knockdowns. The smallest $q$ was
0.18 (GATA1 at the 0.5 cut-off, under the other-knockdown null, where
the smallest possible $p$ is 0.059). We also ran a fifth null. It uses
random groups of control cells of the knockdown's size and runs the
feature selection again on each group. Unlike the third null, it does
not keep as many responding features as the knockdown. Every factor
with at least 2 targets among a responding feature's top 20 genes passed
that null (49 tests, 16 factors). This happens because random control
groups rarely select any feature, and never at cut-offs of 0.25 or
above. So that null shows only that the knockdown changed something, not
that the features know the factor's targets (S3~Appendix).

GATA1 is the factor the deployment singled out. At the cut-off the
deployment used (a difference of more than 0.5 between knockdown and
control cells in a feature's mean activation), seven features responded
to the knockdown. The best had 4 of its 20 top genes among the
221 GATA1 ChIP-seq targets in the gene set, where 3.1 are expected by
chance ($p = 0.32$ under the count-matched null, $0.43$ when gene length
is also matched, $0.19$ against random control groups with the same
number of features); GATA1 ranked 82nd of 291 factors for these
features. The test detects strong but not weak specificity. We planted
GATA1 targets among the top 20 genes of one responding feature, with 200
simulations for each number of planted targets. With 6 or more planted
targets, the count-matched test gave $p < 0.05$ in all 200 simulations.
With 5 it did so in 80\% of simulations, and with 4 in 29\%
(Fig~\ref{fig:caseReg}C). These power values count every overlapping
target. The deployed statistic stopped counting at 5. For GATA1 that
version gives the same $p$ here (0.32), but its smallest possible $p$ is
0.099, so it could not detect any planted signal at this cut-off
(S3~Appendix). The data thus exclude a GATA1 feature as specific as the
one the deployment reported (8 of 20), but not a weaker one.

\subsubsection*{Traced circuits do not predict the direction of knockdown responses}
Ablating 120 source features (30 at each of layers 0, 3, 6 and 9) in 200
K562 control cells and reading the change in every later feature gave
318,387 edges between features. A source--target pair counted as an
edge when the target's change had a Cohen's $d$ (its mean change over the
200 cells divided by the standard deviation of that change) above 0.5 in
absolute value, and went the same way in more than 70\% of cells. In
57.9\% of edges, removing the source lowered the target; in the rest it
raised it. We turned these edges into predictions about genes. Each
feature has 10 top genes, the genes on which it is most active. Each edge
pairs every top gene of its source feature with every top gene of its
target feature. We kept a gene pair if at least 2 edges linked it, or one
edge with $|d| > 2$. We predicted that silencing the first gene lowers the
second if more than half of the edges linking them were ones in which
removing the source lowered the target, and that it raises the second
otherwise. We checked these predictions in K562 cells in which the first
gene was silenced by CRISPR interference. The observed direction is the
sign of the second gene's mean $\log(1+\text{CP10k})$ value in the
silenced cells minus its mean in 3,000 control cells. This gave 796,304
records from 248 silenced genes; each record is one pair of a silenced
gene and a measured target gene. The prediction was right
50.3\% of the time (95\% interval over silenced genes 49.9--50.7\%),
below the 52.5\% obtained by always predicting ``down''
(Fig~\ref{fig:caseReg}D). Balanced accuracy, which is 50\% for any
constant rule, was 50.2\% ($+0.19$ points, $-0.18$ to $+0.57$), and the
Matthews correlation was $+0.004$ ($-0.004$ to $+0.011$). A rule that
uses no model, predicting for each target the direction in which it
usually moves when other genes are silenced (fitted on other silenced
genes), was right 58.5\% of the time, 8.3 points more than the circuit
(7.3--9.2). An exploratory split by the size of the observed change found
a balanced accuracy of 51.8\% for changes of 0.1 to 0.25 log-fold, but the
no-model rule reached 68.4\% on the same records.

\subsubsection*{Removing three features at once had an additive effect in the 512 triplets tested}
In 512 triplets of features (eight each at layers 0, 5 and 9) ablated
together and separately in 20 K562 control cells, no triplet had a
three-way interaction at any of the 4,928 last-layer features after BH
correction (2,523,136 tests). The 24 features were drawn at random before
any ablation: at each layer, one from each eighth of activation
frequency, among the features active in at least 1\% of tokens (1,750,
460 and 618 candidates at layers 0, 5 and 9). No annotation was used to
choose them. After subtracting the part expected from cell-to-cell
noise, the share of the joint effect's energy (its squared size) that
was not explained by the three single effects and their pairwise
interactions was 0.0003\% (median over the 512 triplets). Its 95\% basic
bootstrap interval, resampling the 20 cells (1,000 resamples), was
$-0.09\%$ to $+0.06\%$. The interval can go below zero because of the
noise subtraction. The test detected a planted interaction of 5\% of each
cell's joint effect (about 0.2\% of its energy) in 64 of 64 triplets.
Among the 192 feature pairs in these triplets, removing both features
changed the last-layer features by less energy than the two single
changes added together in 5 pairs, and by more in 12 (95\% interval
excluding equality). About 5 per side would be expected by chance, so
these are weak signs. The two clearest pairs both link a layer-5 feature to a layer-9 feature,
with many targets significant after BH correction. In both, removing the
two features together changed the last-layer features by 3--4\% less
energy than the two single removals added together. This fits the
layer-5 feature driving the layer-9 one. Removing all three features together moved the
last-layer residual stream by about 2\% of its norm (median 2.2\%, range
1.6--4.1\%, over 16 cell--triplet pairs), so this describes the model's
response to small edits. The change from single removals was not
recorded. The result covers 8 of the 4,928 features at each of three of
the twelve layers, one read-out layer (layer~11), one cell type and one
SAE training seed. It does not show that three-way interactions are
absent elsewhere in the model.

\subsubsection*{Features that track erythroid maturity move the model's output no further than random features, apart from one feature that marks erythroid cells}
Steering here means scaling one sparse-autoencoder feature up or down
inside the model while it reads a cell, and measuring how far the model's
output moves along a maturity score. The deployment's steering effects
came from the hook that deleted a block, not from the features (see
above). We fixed the cells, score, features, null and tests before any
forward pass. The maturity axis is the erythroid branch of blood
development in Tabula Sapiens bone marrow: hematopoietic stem cells
(HSCs), common myeloid progenitors, erythroid progenitors and
erythrocytes, as labelled in the cell-type annotation. For each cell,
the model's output here is its next-gene scores (logits: one score for
each of the 20,275 tokens in the vocabulary), averaged over all
positions of the cell's sequence. The maturity score is the cosine
similarity of this output with the mean output of erythrocytes, minus
its cosine similarity with the mean output of HSCs. The two means came
from 90 erythrocytes and 76 HSCs of three donors. In cells of two other
donors the score separated erythroid progenitors from HSCs with an AUROC
of 0.991, and its Spearman correlation with stage was 0.84 (95\% bootstrap interval over
cells 0.78--0.88). We give each effect as a share of the span, the mean
score of erythrocytes minus that of HSCs (0.198). The score describes the
model's output; no change in a cell's real state was measured.

We steered 50 HSCs from the two held-out donors (47 of them from one
donor). At each of layers 0, 3, 6, 9 and 11, the three features whose
activity tracked stage most closely were the candidates. Each was
multiplied by 0, 2 or 5 and compared with 20 random features of the same
layer with similar activity. Fourteen of the 15 candidates moved cells no
further in their predicted direction than the random features
(directional $p$ 0.19--1.00; 0.0004\% to 1.5\% of the span at five-fold
strength). The exception, feature 2903 at layer~9, beat all 20 random
features at every strength ($p = 1/21$, the smallest possible; about 0.7
of 15 candidates would reach this by chance) and moved HSCs by 3.5\% of
the span at five-fold strength (95\% bootstrap interval over cells
1.9--5.2\%). It is active at 99\% of gene positions in erythroid
progenitors and at 11\% in HSCs, so it marks the erythroid state. Per
unit of edit it moved cells 37\% as far as a size-matched edit along the
mean difference between erythrocyte and HSC hidden states, and the genes
it raised most (AHSP, ALAS2, HBM, GYPA, HBG1, SLC4A1) are among the ten
that this mean-difference edit raised most. These three checks (where
the feature is active, the comparison per unit of edit and the
most-raised genes) were not planned. Two tests fixed in advance looked at
each layer's three candidates as a group. Neither found them stronger
than random features. First, at no layer did the three candidates push
cells further in their predicted direction than the 1,771 possible sets
of three of that layer's 23 tested features (directional $p$ 0.09--0.64
over the five layers and three strengths). Second, summed over the five
layers, the candidates ranked no higher among the 23 tested features of
their layer than chance (one-sided permutation $p$ 0.24--0.36 over the
three strengths; 100,000 permutations within layers). A third test, also
pooled over layers, found a weak trend. Within each layer, features whose activity rose
more with stage pushed cells a little further toward erythrocytes when
amplified, and a little further back when removed. The mean of the five
within-layer Spearman correlations, each over 23 features, was
0.18--0.19 in size and had the predicted sign (one-sided permutation $p$
0.022--0.031; two-sided 0.045--0.063). Twenty of each layer's 23
features are random features, so this trend does not single out the
candidates. These sparse autoencoders were trained on K562 cells and explained only
34--59\% of the hidden-state variance of the steered HSCs. So single
features of these autoencoders are not maturity switches. This does not
show that the model lacks a maturity direction: adding the full mean
difference moved all cells toward erythrocytes. Other lineages, feature
combinations, autoencoders trained on bone-marrow cells and MaxToki-1B
were not tested (S3~Appendix).

\subsubsection*{Hidden states show no blood developmental order beyond cell-type identity}
The deployed analysis fitted a small read-out head that maps the hidden
states of cell groups (anchors: cells that share donor, tissue, cell type
and stage label) onto a curated tree of 34 blood developmental stages,
and reported a developmental manifold that transfers to new donors. It
judged the head with four pass marks. The first is trustworthiness, at
least 0.80. It checks that anchors that are near neighbours in the
head's output were also close in the head's input, and it does not use
the stage tree. The other three are Spearman correlations between
output distances and stage-tree distances, each at least 0.20. They use
anchors the head was not fitted on, held out at random, by donor or by
lineage branch. The last of these is the branch mark. MaxToki passed
all four on the fitting panel (290 anchors from 11
donors) and on the external test panel (600 anchors from 13 other
donors). The deployment also built a second test panel, which it called
zero-shot: 160 other anchors whose 12 donors are all external-panel
donors, so the two test panels are not independent. On this panel
MaxToki failed the branch mark ($0.147 \pm 0.048$, mean $\pm$ SD over 20
head seeds; 3 of 20 seeds pass). This failure depends on the anchor
set. Some anchors hold only Smart-seq2 cells (Smart-seq2 is a
plate-based sequencing method). When these anchors are dropped (43
fitting-panel, 13 external and 8 zero-shot anchors) and the head is
fitted again, the zero-shot branch mark is $0.34 \pm 0.01$ (mean $\pm$
SD over 10 head seeds; 10 of 10 seeds pass). A random 64-number code for
each cell-type label, which knows nothing about development, passed the three
correlation marks on all three panels. The branch mark can be high just
because anchors with the same label sit together, and only the T lineage
has graded stages. Against a structured null that shuffles stages
between cell-type classes within each branch, MaxToki was not above
the null on the branch mark on any of the three panels. This held both
over all anchor pairs and over only pairs of anchors with different
cell types. MaxToki was also not above the null on cross-lineage
(global) order, the Spearman correlation over all anchor pairs of a test
panel. Across these numbers the one-sided $p$ ran from 0.15 to 0.76 (20
to 2,000 null draws). On pairs of anchors with different cell types, the
branch mark was 0.067, 0.006 and 0.016 on the fitting, external and
zero-shot panels. That is close to zero and well below the
0.20 pass mark. Held-out order within the T lineage was 0.038.

MaxToki was ahead of simple features only in global order
(Fig~\ref{fig:caseGeom}D): 0.865 on the external panel and 0.867 on the
zero-shot panel. That is above the cell-type label code by 0.064 (95\%
donor bootstrap interval 0.050--0.091) and 0.091 (0.027--0.163), and
above expression of highly variable genes by 0.102 (0.073--0.140) and
0.110 (0.038--0.151). A bag of input tokens uses no model. It describes
each cell group by the genes in its cells' model input, and genes nearer
the front of the input get more weight. MaxToki's lead over this bag was
small: 0.016 (0.002--0.046) on the external panel and 0.048
(0.009--0.149) on the zero-shot panel. Over 10 head seeds the lead was
$0.004 \pm 0.005$ on the external panel, a tie, and $0.053 \pm 0.005$ on
the zero-shot panel (above 0 in 10 of 10 seeds). A second bag, built the
same way from the gene order that the deployed analysis fed the model,
also uses no model, and it matched MaxToki on both panels: MaxToki minus
this bag was $-0.010$ ($-0.023$ to 0.016) and $-0.012$ ($-0.031$ to
0.022) (donor bootstrap, 1,000 replicates). So whether MaxToki leads
depends on which gene order the bag uses. Most of global order
reflects lineage and cell-type grouping, which the structured null keeps
(null means 0.843 and 0.831, $p = 0.29$ and 0.15). S3~Appendix gives the
remaining checks, including a negative control that could not fail and
the single attention head singled out by the deployment, which beat
random projections only in trustworthiness.

\begin{figure}[H]
\caption{\textbf{MaxToki-217M: representation geometry after correction.}
(A)~Input gene-embedding tables: Pearson correlation between MaxToki's
gene--gene cosine-similarity matrix and that of Geneformer~V2-316M or
scGPT, on three gene panels from three tissue datasets (350--382 genes;
line, 95\% gene bootstrap interval). Dashed line, chance (genes
shuffled). (B)~Effective rank of the per-gene mean vectors at each layer,
for the trained model and three randomly initialised copies of the same
architecture, on the same 2,000 Tabula Sapiens immune cells and 1,500
genes (band, 95\% jackknife interval over 10 blocks of 200 cells).
(C)~Similarity (linear CKA) of the per-gene mean vectors between three
disjoint 2,000-cell samples, from layer~1, for the trained model and one
randomly initialised copy (band, 95\% jackknife interval over 30 groups
of 50 genes). (D)~Cross-lineage developmental order of cell-group
centroids on the external and zero-shot panels, for MaxToki-217M and for
three representations that use no model: the bag of each cell's input
tokens, a lookup of cell-type labels, and highly variable gene
expression (line, 95\% donor bootstrap interval, 2,000 replicates).}
\label{fig:caseGeom}
\end{figure}

\subsubsection*{Effective rank falls with depth in randomly initialised copies too}
We passed three disjoint samples of 2,000 Tabula Sapiens immune
cells~\cite{tabulasapiens2022} through MaxToki-217M and, at each layer,
averaged the hidden state of each of 1,500 genes over the cells that
contain it. Effective rank~\cite{roy2007effective} counts how many
directions these gene vectors use (at most 1,232, the hidden size). In the
trained model it was 561 at the token embedding (layer~0), 328 after the
first block and 227 at layer~11 (95\% jackknife interval over ten
200-cell blocks 224.8--229.5; 225.6 and 226.3 in the other two samples)
(Fig~\ref{fig:caseGeom}B). The deployed analysis, whose inputs were
encoded on the wrong scale, reported 94 at layer~11 for the same cells
and genes. Three randomly initialised copies of the architecture fell
from 816--817 at layer~0 to 390--421 at layer~11. From layer~1 to
layer~11 every random copy lost more of its effective rank than the
trained model (34.6--39.6\% against 30.8--31.3\%; the cell-block
jackknife intervals do not overlap), so the fall with depth does not need
training. Training lowers the effective rank at every layer, starting at
the embedding table, and adds two drops of its own: in the first block
(41.5\% against 21.0--21.3\% with random weights) and in the last step,
from layer~10 to layer~11, which runs the last block and then the final
normalisation (14.6--15.1\%, against a rise of 1.0--2.1\%; these two
parts were not measured separately). A feature-shuffle null cannot
separate these cases, because it stays far above the real effective rank
in every model.

Linear centred kernel alignment (CKA)~\cite{kornblith2019cka} measures how alike two sets of gene
vectors are, up to rotation and scale (1 means the same geometry).
Between the three samples it was 0.9997 at layer~1 and 0.9962 at
layer~11 in the trained model (95\% jackknife interval over 30 groups of
50 genes 0.9959--0.9964), against 0.9938 and 0.9647 (0.9598--0.9697) in a
randomly initialised copy (Fig~\ref{fig:caseGeom}C). In the trained model
most of the gap from 1 is sampling noise in the per-gene averages: with
1,000 instead of 2,000 cells per sample, $1-\mathrm{CKA}$ at layer~11
rose from 0.0038 to 0.0073. The trained model's gene averages are
therefore more stable across cell samples than a random model's, but
these values mainly measure how far a gene's average moves between random
sets of cells, and they should not be compared with CKA values reported
for other models. Other tissues, other initialisation schemes and
MaxToki-1B were not tested (S3~Appendix).

\subsubsection*{Cross-model alignment of gene embeddings}
This analysis compares the input gene-embedding tables of the three
models, which do not depend on how cells are encoded. We used three
overlapping panels of 350 to 382 genes, taken from three tissue
datasets: Tabula Sapiens lung, Tabula Sapiens immune cells, and an
external lung dataset. For each model we computed the cosine similarity
of every gene pair. We then took the Pearson correlation between two
models' similarity matrices. Between MaxToki-217M and scGPT this
correlation was 0.26--0.28. When genes were shuffled it was 0.00 (mean;
SD 0.004). We also report a held-out canonical correlation. Canonical
correlation measures how well linear projections of the two tables line
up. The held-out version is fitted on 80\% of the genes and scored on
the other 20\%, averaged over 10 random splits. Between MaxToki-217M and
scGPT it was 0.481--0.495, with chance about 0.00. MaxToki-217M agreed
more with Geneformer~V2-316M, which shares its gene vocabulary and token
identifiers: Pearson 0.39--0.40 and held-out canonical correlation
0.558--0.604 (Fig~\ref{fig:caseGeom}A).
The Geneformer-minus-scGPT difference in Pearson correlation was 0.10 to
0.14, with gene-bootstrap intervals that exclude zero on every panel. The
in-sample canonical correlation (fitted and scored on the same genes),
often used for such comparisons, is high even for unrelated tables. In
the deployed setting (each table reduced to 30 principal components, mean
of 10 canonical correlations) it was 0.41--0.43 for shuffled genes,
against 0.77--0.78 for Geneformer and 0.72--0.74 for scGPT. Its chance
level depends on the number of genes and of components. In the setting
the Study~A executors used (20 components, mean of all 20), chance was
about 0.20. So we report the held-out value. These are input tables only; they say nothing
about the models' internal representations.

